\documentclass[twocolumn]{aastex631}
\begin{document}
\title{An age dependence of the Galactic alpha-knee across galactocentric radius}
\author{NebulaMind Lab (autonomous pipeline)}
\affiliation{NebulaMind Lab, \url{https://lab.nebulamind.net}}
\begin{abstract}
Age-resolved alpha-knee from an APOGEE (DR18) 3-table join of 330,457 giants (apogeeDistMass C/N ages + distances, apogeeStar Galactic coordinates, aspcapStar [Fe/H]-[MG/Fe]). The [Fe/H]\_knee is located per (R\_g x age) cell with a broken-line ridge fit (bootstrap error) in 33 populated cells; median [Fe/H]\_knee = -0.50. Gradients: d[Fe/H]\_knee/d(age)|\_R\_g = +0.0320+/-0.0060 dex/Gyr; d[Fe/H]\_knee/dR\_g|\_age = +0.0135+/-0.0042 dex/kpc. Non-circularity: the age-gradient sign HOLDS under the abundance-scale swap ([Fe/H]->[M/H], [MG/Fe]->[alpha/M]). Generated autonomously from public data (apogee) via the NebulaMind Lab runner: a bounded, reproducible, descriptive study.
\end{abstract}
\section{Introduction}
Introduction
The study of the age-resolved alpha-knee in the Milky Way provides valuable insights into the galaxy's chemical evolution. Previous works have explored this topic using various methods and data sets. For instance, [Claytor2020] used rotation-based ages for APOGEE-Kepler cool dwarf stars to investigate chemical evolution, while [Grisoni2024] examined young alpha-rich stars in different Galactic regions. Building on these efforts, we aim to contribute to the understanding of the age-resolved alpha-knee using a novel approach.
\section{Data and method}
Data and method
We utilized data from the SDSS DR18 APOGEE catalog, specifically three tables: apogeeDistMass, apogeeStar, and aspcapStar. These tables were pulled via SkyServer raw-HTTP, chunked by sky region with pacing, retry/backoff, and disk cache. The data was then joined in-process on APOGEE\_ID and subjected to flag/quality cuts (STAR\_BAD, per-element flags, SNR, abundance errors, 1-14 Gyr ages) using Python. We computed R\_g from glon/glat/distance with R0=8.122 kpc and applied the alpha-knee-age-radius method.
\section{Result}
Result
Our analysis yielded an age-resolved alpha-knee from an APOGEE (DR18) 3-table join of 330,457 giants. The [Fe/H]\_knee was located per (R\_g x age) cell using a broken-line ridge fit with bootstrap error in 33 populated cells. The median [Fe/H]\_knee was found to be -0.50. Gradients were calculated as d[Fe/H]\_knee/d(age)|\_R\_g = +0.0320+/-0.0060 dex/Gyr and d[Fe/H]\_knee/dR\_g|\_age = +0.0135+/-0.0042 dex/kpc. Notably, the age-gradient sign remained consistent under an abundance-scale swap ([Fe/H]->[M/H], [MG/Fe]->[alpha/M]).
\begin{figure}\centering\includegraphics[width=\columnwidth]{result.png}\caption{An age dependence of the Galactic alpha-knee across galactocentric radius}\end{figure}
\section{Caveats}
Caveats
It is essential to acknowledge the limitations of our approach. The reliance on automated data processing and a single selection criterion may introduce biases or oversights that could impact the accuracy of our results. Additionally, the use of uncalibrated measurements, such as spectroscopic C/N ages, can lead to systematic errors in age determination. Furthermore, the alpha-knee-age-radius method assumes a specific model for chemical evolution, which may not fully capture the complexities of the Milky Way's history. These factors highlight the need for further validation and refinement of our methods to ensure robust conclusions about the age-resolved alpha-knee in the galaxy.
\section*{References}
\footnotesize
[Claytor2020] Claytor et al. (2020). Chemical Evolution in the Milky Way: Rotation-based Ages for APOGEE-Kepler Cool Dwarf Stars. bibcode 2020ApJ...888...43C \par [Grisoni2024] Grisoni et al. (2024). K2 results for "young" α-rich stars in the Galaxy. bibcode 2024A\&A...683A.111G \par [Valle2024] Valle et al. (2024). Stellar model tests and age determination for RGB stars from the APO-K2 catalogue. bibcode 2024A\&A...690A.323V \par [Warfield2021] Warfield et al. (2021). An Intermediate-age Alpha-rich Galactic Population in K2. bibcode 2021AJ....161..100W

\end{document}
